% ============================================================================
% CHAPTER 9: INSTRUMENTAL VARIABLES
% ============================================================================

\chapter{Instrumental Variables}
\label{ch:iv}

% ============================================================================
\section{Introduction}
\label{sec:iv_intro}
% ============================================================================

When treatment assignment is correlated with unobserved determinants of the outcome, ordinary least squares estimates are biased. \emph{Instrumental variables} (IV) estimation provides a solution when we can find a variable that affects treatment but has no direct effect on the outcome.

\subsection{The Endogeneity Problem}

Consider the structural equation:
\begin{equation}
Y_i = \beta D_i + X_i'\gamma + \varepsilon_i
\label{eq:structural}
\end{equation}
where $D_i$ is the endogenous treatment and $\varepsilon_i$ is the error term. Endogeneity arises when:
\begin{equation}
\Cov(D_i, \varepsilon_i) \neq 0
\end{equation}

This correlation can arise from:
\begin{enumerate}
    \item \textbf{Omitted variables}: Unobserved confounders affect both $D$ and $Y$
    \item \textbf{Simultaneity}: $Y$ affects $D$ (reverse causality)
    \item \textbf{Measurement error}: $D$ is measured with error
\end{enumerate}

\subsection{The IV Solution}

An \emph{instrument} $Z$ provides exogenous variation in treatment. The key idea: use only the variation in $D$ that is driven by $Z$, which is uncorrelated with $\varepsilon$.

\begin{definition}[Instrumental Variable]
\label{def:instrument}
A variable $Z$ is a valid instrument for the effect of $D$ on $Y$ if:
\begin{enumerate}
    \item \textbf{Relevance}: $\Cov(Z, D) \neq 0$ (testable)
    \item \textbf{Exclusion}: $\Cov(Z, \varepsilon) = 0$ (maintained assumption)
\end{enumerate}
\end{definition}

\subsection{Historical Examples}

Classic instrumental variables include:
\begin{itemize}
    \item \textbf{Draft lottery}: Random draft numbers as instruments for military service
    \item \textbf{Quarter of birth}: Birth quarter as instrument for education \cite{angrist1991draft}
    \item \textbf{Geographic proximity}: Distance to college as instrument for education
\end{itemize}

\subsection{Chapter Roadmap}

This chapter develops:
\begin{itemize}
    \item Section~\ref{sec:iv_identification}: Formal identification conditions
    \item Section~\ref{sec:iv_wald}: The Wald estimator for just-identified models
    \item Section~\ref{sec:iv_2sls}: Two-stage least squares estimation
    \item Section~\ref{sec:iv_liml}: LIML and Fuller estimators
    \item Section~\ref{sec:iv_gmm}: GMM estimation with overidentification
    \item Section~\ref{sec:iv_inference}: Inference and standard errors
    \item Section~\ref{sec:iv_implementation}: Python implementation
    \item Section~\ref{sec:iv_validation}: Monte Carlo validation
\end{itemize}


% ============================================================================
\section{Identification}
\label{sec:iv_identification}
% ============================================================================

\subsection{The Structural and Reduced Form Equations}

The complete IV model consists of:

\textbf{Structural equation} (of interest):
\begin{equation}
Y_i = \beta D_i + X_i'\gamma + \varepsilon_i
\end{equation}

\textbf{First stage} (treatment equation):
\begin{equation}
D_i = Z_i'\pi + X_i'\delta + \nu_i
\label{eq:first_stage}
\end{equation}

\textbf{Reduced form} (outcome on instruments):
\begin{equation}
Y_i = Z_i'\rho + X_i'\theta + \eta_i
\label{eq:reduced_form}
\end{equation}

\subsection{Identification Conditions}

\begin{assumption}[Relevance]
\label{ass:relevance}
The instruments are correlated with the endogenous variable:
\begin{equation}
\E[Z_i D_i'] \neq 0
\end{equation}
Equivalently, $\pi \neq 0$ in the first stage equation.
\end{assumption}

\begin{assumption}[Exclusion Restriction]
\label{ass:exclusion}
The instruments are uncorrelated with the structural error:
\begin{equation}
\E[Z_i \varepsilon_i] = 0
\end{equation}
\end{assumption}

\begin{assumption}[Rank Condition]
\label{ass:rank}
For $q$ instruments and $p$ endogenous variables:
\begin{equation}
\text{rank}(\E[Z_i D_i']) = p
\end{equation}
\end{assumption}

\begin{proposition}[Order Condition]
A necessary condition for identification is $q \geq p$: at least as many instruments as endogenous variables.
\end{proposition}

\subsection{DAG Representation}

The IV assumptions can be represented graphically:
\begin{itemize}
    \item $Z \rightarrow D$: Relevance (instrument affects treatment)
    \item $Z \not\rightarrow Y$: Exclusion (no direct effect on outcome)
    \item $Z \not\leftarrow U \rightarrow Y$: No confounding of $Z$-$Y$ relationship
\end{itemize}

where $U$ represents unobserved confounders.


% ============================================================================
\section{The Wald Estimator}
\label{sec:iv_wald}
% ============================================================================

In the simplest case with one instrument and one endogenous variable (just-identified), the IV estimator takes a simple form.

\subsection{Binary Instrument Case}

With a binary instrument $Z \in \{0, 1\}$:

\begin{theorem}[Wald Estimator]
\label{thm:wald}
Under assumptions of relevance and exclusion, the causal effect is:
\begin{equation}
\beta = \frac{\E[Y | Z=1] - \E[Y | Z=0]}{\E[D | Z=1] - \E[D | Z=0]} = \frac{\text{Reduced form}}{\text{First stage}}
\label{eq:wald}
\end{equation}
\end{theorem}

\begin{proof}
From the structural equation $Y = \beta D + \varepsilon$ and exclusion $\E[\varepsilon | Z] = 0$:
\begin{align}
\E[Y | Z=1] - \E[Y | Z=0] &= \beta(\E[D | Z=1] - \E[D | Z=0])
\end{align}
Rearranging gives the result.
\end{proof}

\subsection{Interpretation}

The Wald estimator can be understood as:
\begin{equation}
\hat{\beta}_{\text{Wald}} = \frac{\hat{\rho}}{\hat{\pi}}
\end{equation}
where $\hat{\rho}$ is the reduced form coefficient and $\hat{\pi}$ is the first stage coefficient.

\begin{insight}
The Wald estimator divides the ``intent-to-treat'' effect (reduced form) by the ``compliance rate'' (first stage). This gives the effect of treatment on those who comply with the instrument.
\end{insight}


% ============================================================================
\section{Two-Stage Least Squares}
\label{sec:iv_2sls}
% ============================================================================

Two-stage least squares (2SLS) generalizes the Wald estimator to multiple instruments and covariates.

\subsection{The 2SLS Procedure}

\begin{definition}[Two-Stage Least Squares]
\label{def:2sls}
The 2SLS estimator is computed in two stages:
\begin{enumerate}
    \item \textbf{First stage}: Regress $D$ on $Z$ and $X$, obtain fitted values $\hat{D}$
    \item \textbf{Second stage}: Regress $Y$ on $\hat{D}$ and $X$
\end{enumerate}
\end{definition}

\subsection{Matrix Formulation}

Let $W = [D, X]$ be the matrix of endogenous and exogenous regressors, and $\tilde{Z} = [Z, X]$ the full instrument matrix. The 2SLS estimator is:

\begin{equation}
\hat{\beta}_{2\text{SLS}} = (W' P_{\tilde{Z}} W)^{-1} W' P_{\tilde{Z}} Y
\label{eq:2sls_matrix}
\end{equation}

where $P_{\tilde{Z}} = \tilde{Z}(\tilde{Z}'\tilde{Z})^{-1}\tilde{Z}'$ is the projection matrix.

\begin{theorem}[Consistency of 2SLS]
\label{thm:2sls_consistency}
Under the rank condition and exclusion restriction:
\begin{equation}
\hat{\beta}_{2\text{SLS}} \convp \beta
\end{equation}
\end{theorem}

\subsection{Variance Estimation}

The asymptotic variance of 2SLS depends on the error structure:

\textbf{Homoskedastic variance}:
\begin{equation}
\hat{V}_{2\text{SLS}} = \hat{\sigma}^2 (W' P_{\tilde{Z}} W)^{-1}
\end{equation}
where $\hat{\sigma}^2 = \frac{1}{n-k}\sum_i \hat{\varepsilon}_i^2$.

\textbf{Robust (heteroskedasticity-consistent) variance}:
\begin{equation}
\hat{V}_{\text{robust}} = (W' P_{\tilde{Z}} W)^{-1} \left(\sum_i \hat{\varepsilon}_i^2 \hat{w}_i \hat{w}_i'\right) (W' P_{\tilde{Z}} W)^{-1}
\end{equation}
where $\hat{w}_i = P_{\tilde{Z}} w_i$.

\subsection{Just-Identified vs. Over-Identified}

\begin{itemize}
    \item \textbf{Just-identified} ($q = p$): Exactly as many instruments as endogenous variables. 2SLS is unique but may be biased in small samples.
    \item \textbf{Over-identified} ($q > p$): More instruments than needed. 2SLS gains efficiency but requires testing overidentifying restrictions.
\end{itemize}


% ============================================================================
\section{LIML and Fuller Estimators}
\label{sec:iv_liml}
% ============================================================================

Limited Information Maximum Likelihood (LIML) and Fuller estimators address the finite-sample bias of 2SLS.

\subsection{k-Class Estimators}

The k-class family of estimators is defined as:
\begin{equation}
\hat{\beta}_k = (W'(I - k M_{\tilde{Z}})W)^{-1} W'(I - k M_{\tilde{Z}})Y
\label{eq:kclass}
\end{equation}
where $M_{\tilde{Z}} = I - P_{\tilde{Z}}$.

Special cases:
\begin{itemize}
    \item $k = 0$: OLS (ignores endogeneity)
    \item $k = 1$: 2SLS
    \item $k = \hat{\lambda}_{\min}$: LIML
\end{itemize}

\subsection{LIML Estimator}

\begin{definition}[LIML]
\label{def:liml}
LIML uses $k = \hat{\lambda}_{\min}$, the smallest eigenvalue of:
\begin{equation}
(W'M_X W)^{-1}(W'M_{[X,Z]} W)
\end{equation}
where $M_X = I - X(X'X)^{-1}X'$.
\end{definition}

\begin{proposition}[LIML Properties]
LIML is:
\begin{enumerate}
    \item Less biased than 2SLS with weak instruments
    \item Has higher variance than 2SLS
    \item Median-unbiased under normality
\end{enumerate}
\end{proposition}

\subsection{Fuller Estimator}

\begin{definition}[Fuller Estimator]
\label{def:fuller}
The Fuller estimator modifies LIML with a bias correction:
\begin{equation}
k_{\text{Fuller}} = \hat{\lambda}_{\min} - \frac{\alpha}{n - L}
\end{equation}
where $L$ is the number of excluded instruments and $\alpha$ is a tuning parameter (typically $\alpha = 1$ or $\alpha = 4$).
\end{definition}

Fuller-1 ($\alpha = 1$) provides a good bias-variance tradeoff across instrument strength regimes.

\subsection{When to Use Each Estimator}

\begin{table}[ht]
\centering
\caption{Estimator Selection Guide}
\begin{tabular}{lcc}
\toprule
\textbf{Scenario} & \textbf{Recommended} & \textbf{Rationale} \\
\midrule
Strong instruments ($F > 20$) & 2SLS & Efficient, minimal bias \\
Moderate ($10 < F < 20$) & Fuller-1 & Balance bias-variance \\
Weak ($F < 10$) & LIML or Fuller-4 & Less biased \\
\bottomrule
\end{tabular}
\end{table}


% ============================================================================
\section{GMM Estimation}
\label{sec:iv_gmm}
% ============================================================================

Generalized Method of Moments (GMM) provides an efficient estimator when there are more instruments than endogenous variables.

\subsection{Moment Conditions}

The exclusion restriction implies moment conditions:
\begin{equation}
\E[Z_i \varepsilon_i] = \E[Z_i(Y_i - D_i'\beta - X_i'\gamma)] = 0
\label{eq:moment_conditions}
\end{equation}

\subsection{GMM Objective}

The GMM estimator minimizes:
\begin{equation}
Q(\beta) = \left(\frac{1}{n}\sum_i Z_i \hat{\varepsilon}_i\right)' W \left(\frac{1}{n}\sum_i Z_i \hat{\varepsilon}_i\right)
\label{eq:gmm_objective}
\end{equation}
where $W$ is a positive definite weighting matrix.

\subsection{One-Step vs. Two-Step GMM}

\textbf{One-step GMM}: Uses $W = (Z'Z)^{-1}$, equivalent to 2SLS.

\textbf{Two-step GMM}:
\begin{enumerate}
    \item Compute one-step GMM residuals $\hat{\varepsilon}^{(1)}$
    \item Form optimal weighting: $\hat{W} = \left(\frac{1}{n}\sum_i \hat{\varepsilon}_i^{(1)2} Z_i Z_i'\right)^{-1}$
    \item Re-estimate with $\hat{W}$
\end{enumerate}

Two-step GMM is efficient under heteroskedasticity.

\subsection{Hansen J-Test}

\begin{definition}[Hansen J-Test]
\label{def:hansen_j}
The overidentification test statistic is:
\begin{equation}
J = n \cdot Q(\hat{\beta}) \sim \chi^2_{q-p}
\end{equation}
under the null that all instruments are valid.
\end{definition}

A significant J-statistic suggests at least one instrument may violate the exclusion restriction.


% ============================================================================
\section{Inference}
\label{sec:iv_inference}
% ============================================================================

\subsection{Standard Errors}

Three variance estimation approaches are available:

\begin{enumerate}
    \item \textbf{Homoskedastic}: Assumes $\Var(\varepsilon | Z) = \sigma^2$
    \item \textbf{Robust}: Allows heteroskedasticity
    \item \textbf{Clustered}: Allows within-cluster correlation
\end{enumerate}

\subsection{First-Stage F-Statistic}

\begin{definition}[First-Stage F-Statistic]
\label{def:first_stage_f}
The F-statistic tests joint significance of excluded instruments in the first stage:
\begin{equation}
F = \frac{(R^2_{\text{full}} - R^2_{\text{reduced}})/(q)}{(1 - R^2_{\text{full}})/(n - k)}
\end{equation}
\end{definition}

\begin{remark}[Rule of Thumb]
Staiger and Stock (1997) suggest $F > 10$ indicates instruments are not weak. This is a rough guideline; see Chapter~\ref{ch:weak_iv} for formal testing.
\end{remark}

\subsection{Confidence Intervals}

Standard Wald confidence intervals:
\begin{equation}
\hat{\beta} \pm z_{\alpha/2} \cdot \hat{\text{SE}}(\hat{\beta})
\end{equation}

These may have poor coverage with weak instruments; see Chapter~\ref{ch:weak_iv} for robust alternatives.


% ============================================================================
\section{Implementation}
\label{sec:iv_implementation}
% ============================================================================

\subsection{Python API}

The \texttt{causal\_inference.iv} module provides IV estimators:

\begin{minted}{python}
from causal_inference.iv import TwoStageLeastSquares, LIML, Fuller, GMM

# Prepare data
Y = outcome_data      # n x 1
D = treatment_data    # n x 1 (endogenous)
Z = instrument_data   # n x q
X = covariate_data    # n x p (optional)

# 2SLS estimation
tsls = TwoStageLeastSquares()
tsls.fit(Y, D, Z, X)
print(f"Estimate: {tsls.coef_[0]:.4f}")
print(f"SE: {tsls.se_[0]:.4f}")
print(f"First-stage F: {tsls.first_stage_f_:.2f}")

# LIML for potentially weak instruments
liml = LIML()
liml.fit(Y, D, Z, X)
print(f"LIML estimate: {liml.coef_[0]:.4f}")
print(f"LIML kappa: {liml.kappa_:.4f}")

# Fuller-1 (balanced bias-variance)
fuller = Fuller(alpha_param=1.0)
fuller.fit(Y, D, Z, X)

# Two-step GMM with J-test
gmm = GMM(steps='two')
gmm.fit(Y, D, Z, X)
print(f"J-statistic: {gmm.j_statistic_:.2f}")
print(f"J p-value: {gmm.j_pvalue_:.4f}")
\end{minted}

\subsection{Robust Standard Errors}

\begin{minted}{python}
# Heteroskedasticity-robust SEs
tsls_robust = TwoStageLeastSquares(inference='robust')
tsls_robust.fit(Y, D, Z, X)

# Cluster-robust SEs
tsls_cluster = TwoStageLeastSquares(
    inference='clustered',
    clusters=cluster_ids
)
tsls_cluster.fit(Y, D, Z, X)
\end{minted}

\subsection{Complete Example}

\begin{minted}{python}
import numpy as np
from causal_inference.iv import TwoStageLeastSquares
from causal_inference.iv.diagnostics import classify_instrument_strength

# Simulate IV data
np.random.seed(42)
n = 1000
Z = np.random.normal(0, 1, (n, 1))  # Instrument
U = np.random.normal(0, 1, n)        # Unobserved confounder
D = 0.5 * Z.flatten() + 0.5 * U + np.random.normal(0, 0.5, n)
Y = 2.0 * D + U + np.random.normal(0, 1, n)

# Estimate
tsls = TwoStageLeastSquares(inference='robust')
tsls.fit(Y, D, Z)

# Results
print(f"True effect: 2.0")
print(f"2SLS estimate: {tsls.coef_[0]:.4f}")
print(f"95% CI: [{tsls.ci_lower_[0]:.4f}, {tsls.ci_upper_[0]:.4f}]")
print(f"First-stage F: {tsls.first_stage_f_:.1f}")

# Classify instrument strength
strength = classify_instrument_strength(tsls.first_stage_f_, n_instruments=1)
print(f"Instrument strength: {strength}")
\end{minted}


% ============================================================================
\section{Monte Carlo Validation}
\label{sec:iv_validation}
% ============================================================================

\subsection{Data Generating Process}

We validate IV estimators using:
\begin{align}
Y_i &= \beta D_i + \varepsilon_i, \quad \beta = 2.0 \\
D_i &= \pi Z_i + \nu_i \\
(\varepsilon_i, \nu_i) &\sim N(0, \Sigma), \quad \Sigma = \begin{pmatrix} 1 & \rho \\ \rho & 1 \end{pmatrix}
\end{align}
where $\rho = 0.5$ creates endogeneity.

\subsection{Validation Results}

\begin{table}[ht]
\centering
\caption{Monte Carlo Results: IV Estimators ($n=500$, 5,000 replications)}
\label{tab:iv_mc}
\begin{tabular}{lccccc}
\toprule
\textbf{Estimator} & \textbf{True} & \textbf{Mean Est.} & \textbf{Bias} & \textbf{RMSE} & \textbf{Coverage} \\
\midrule
OLS (biased) & 2.0 & 2.50 & 0.50 & 0.51 & 0\% \\
2SLS & 2.0 & 2.01 & 0.01 & 0.12 & 94.8\% \\
LIML & 2.0 & 2.00 & 0.00 & 0.13 & 95.2\% \\
Fuller-1 & 2.0 & 2.00 & 0.00 & 0.12 & 95.0\% \\
GMM (2-step) & 2.0 & 2.01 & 0.01 & 0.12 & 94.9\% \\
\bottomrule
\end{tabular}
\end{table}

Key findings:
\begin{itemize}
    \item OLS is severely biased due to endogeneity
    \item All IV estimators recover the true effect
    \item Coverage rates are near nominal 95\%
    \item 2SLS, Fuller-1, and GMM have similar efficiency
\end{itemize}


% ============================================================================
\section{Practical Guidance}
\label{sec:iv_guidance}
% ============================================================================

\subsection{Assessing Instrument Validity}

\begin{enumerate}
    \item \textbf{Relevance}: Check first-stage F-statistic ($F > 10$ rule of thumb)
    \item \textbf{Exclusion}: Cannot be tested directly---requires economic reasoning
    \item \textbf{Monotonicity} (for LATE): Treatment response is monotone in instrument
\end{enumerate}

\subsection{Common Mistakes}

\begin{itemize}
    \item \textbf{Weak instruments}: Using instruments with low predictive power leads to biased estimates
    \item \textbf{Invalid instruments}: Direct effect of $Z$ on $Y$ violates exclusion
    \item \textbf{Heterogeneous effects}: IV estimates LATE, not ATE, when effects vary
\end{itemize}

\subsection{Reporting Recommendations}

Always report:
\begin{enumerate}
    \item First-stage F-statistic
    \item Both 2SLS and at least one alternative (LIML or Fuller)
    \item Robust standard errors
    \item Hansen J-test if over-identified
\end{enumerate}


% ============================================================================
\section{Summary}
\label{sec:iv_summary}
% ============================================================================

\textbf{Key Results:}
\begin{enumerate}
    \item IV estimation requires relevance (testable) and exclusion (assumed)
    \item 2SLS is the standard estimator for strong instruments
    \item LIML and Fuller reduce bias when instruments are weak
    \item GMM is efficient with heteroskedasticity and multiple instruments
    \item Hansen J-test checks overidentifying restrictions
\end{enumerate}

\textbf{Key References:}
\begin{itemize}
    \item Angrist \& Krueger (1991): Quarter of birth instruments
    \item Stock \& Yogo (2005): Weak instrument testing
    \item Bound, Jaeger \& Baker (1995): Weak instrument bias
\end{itemize}

\textbf{Next Chapter:} Chapter~\ref{ch:weak_iv} develops diagnostics and robust inference for weak instruments.
